\documentclass[10pt]{article}
\usepackage{amsmath,amssymb,geometry}
\geometry{margin=1in}
\title{Four generators for six atoms: a counterexample to 00000008536}
\author{Local AI-assisted solution draft}
\date{3 October 2026}
\usepackage{url}
% Compact consistent title for a short mathematical note.
\makeatletter
\renewcommand{\maketitle}{\begin{center}
{\Large\bfseries\@title\par}\vspace{0.6em}
{\small\@author\par\@date}
\end{center}\vspace{0.5em}}
\makeatother
\setlength{\emergencystretch}{3em}
\begin{document}
\maketitle
\paragraph{Claim being disproved.}
The conjecture asserts that the least number of generators of a distributive lattice equals its number of join-irreducible elements. We use lattice generation under binary meet and join, even without bottom or top as available constants. Thus the example also works under the convention that constants are available.
\paragraph{Construction.}
Let $E=\{01,02,03,12,13,23\}$ be the six edges of the complete graph on vertices $0,1,2,3$. Its Boolean lattice $L=\mathcal P(E)$ has $64$ elements, meet given by intersection, and join given by union. It is distributive. For each vertex $i$, let $S_i$ be its incident edge set:
\[
\begin{aligned}
 S_0&=\{01,02,03\}, & S_1&=\{01,12,13\},\\
 S_2&=\{02,12,23\}, & S_3&=\{03,13,23\}.
\end{aligned}
\]
\paragraph{Generation.}
For distinct $i,j$,
\[
 S_i\cap S_j=\{ij\}.
\]
Therefore every singleton edge set is generated by the four stars. Also $S_0\cap S_1\cap S_2=\varnothing$, so even the bottom is generated without a constant symbol. Taking unions of the appropriate singleton sets generates every member of $\mathcal P(E)$, including $E$ itself. Hence the minimum number of lattice generators is at most $4$.
\paragraph{Join-irreducibles.}
In a Boolean lattice, the nonzero join-irreducibles are exactly the singletons. A singleton cannot be the union of two proper subsets. Conversely, a set of at least two elements is the union of a singleton and its nonempty proper complement inside that set, so it is join-reducible. Thus $L$ has precisely $6$ join-irreducibles. The proposed equality would give minimum generator number $6$, contradicting the explicit generating family of size $4$. No claim about the exact minimum is required for this contradiction.
\paragraph{Formal verification and semantic correspondence.}
The accompanying \texttt{Main.lean} uses only Lean 4.19.0 and \texttt{Std}. The $64$ subsets are represented by six-bit masks. Meet and join are bitwise intersection and union. Their coordinate descriptions and injectivity of the bit representation are checked, and the distributive-lattice identities are proved from the Boolean identities at each coordinate. The bottom property is also checked.

\texttt{JoinIrreducible} is the usual condition $x\ne0$ and $x=a\vee b\Longrightarrow x=a$ or $x=b$. Both the exact six-element list and its cardinality are verified. The inductive type \texttt{Term} has only variables, binary meet, and binary join. The four stars are masks $7,25,42,52$. Explicit terms generate the empty set and each singleton; \texttt{canonicalTerm} takes their appropriate joins to represent any target subset. The theorem \path{four_generators_suffice} proves that every lattice element is the evaluation of a term in those four generators.

If the conjectured minimum equals the join-irreducible count, then every generating $k$-tuple has size at least that count. The proposition \path{ClaimedGeneratorLowerBound} expresses this necessary consequence for the displayed distributive lattice, and \texttt{conjecture8536\_false} directly negates it using the four verified generators and the verified count $6$. Thus the formal contradiction includes actual lattice generation and actual join-irreducibility, rather than relying on an externally asserted interpretation of two numbers.

All finite checks use kernel reduction through \texttt{decide}. The printed axiom lists contain only the standard axiom \texttt{propext}; there are no proof placeholders, external decision oracles, or additional axioms.
\paragraph{Reproduction.}
Run \texttt{lean Main.lean} with Lean 4.19.0.
\end{document}
